\begin{abstract}
Shadow-feature selection keeps a feature when its importance beats that of artificial noise columns, usually row-permuted copies of the real features; Boruta is the best-known version. With a rare positive class this recipe drops features that do predict the target. We show that what controls the loss is the number of minority-class samples $\nmin$ combined with the effect size, and not the imbalance ratio. For the first tree of a gradient-boosted ensemble we prove that the split gain is Pearson's chi-square statistic, that as the prevalence goes to zero with $\nmin$ fixed the root split depends on the data only through $\nmin$ and the placement of the positives among the negatives, and that on a fixed range of thresholds its power depends only on $s = \nmin\delta^2$ for a standardised mean difference $\delta$. The false positive level of the max-shadow rule is $1/(K+1)$ for $K$ shadows at any imbalance, so imbalance takes away power and leaves the level unchanged. A rank test built from the same counting process has power that depends only on a rank analogue of $s$ for any continuous feature, and is more efficient than any single split for Gaussian-like features. We also show that the default intercept of XGBoost 1.x and 2.x adds a toll of order $\lambda/\hat p$ to every split. Simulations against acceptance criteria fixed in advance agree with every proved statement, and full ensembles fail at least as often as the root-split theory predicts. Across fourteen selection methods, eight departures from the Gaussian model and five real tables, a Mann--Whitney screen loses the fewest relevant features at the same false positive rate when the signal shifts the class distributions, and misses signals carried only by spread or by interactions, which trees find. We give guidance in terms of $\nmin$ and a scikit-learn compatible package.
\end{abstract}
